We study the complexity of learning a single reusable abstraction (helper)
that maximizes compression utility over a corpus of programs. We show that
even a highly restricted version of this problem is NP-complete.

\par\noindent
\paragraph{High-level intuition.}
The restricted setting removes almost all sources of complexity usually
associated with library learning. Programs have no internal structure:
each program is just a flat tuple of symbols, and a helper can only specify
that certain coordinates must take fixed values. Thus, learning a helper
reduces to selecting a subset of positions that agree across many programs.

Despite this extreme simplification, the problem remains hard. The reason is
that the learner must jointly choose (i) which positions to constrain, and
(ii) which programs will satisfy those constraints. These two choices
interact combinatorially: constraining more positions makes matches rarer,
while relaxing constraints increases coverage. This trade-off mirrors dense
substructure problems such as biclique detection.

The general version of library learning adds further complexity. Programs
are no longer flat tuples but trees, so a helper must identify a reusable
subtree rather than fixed positions. Variables (or \emph{holes}) allow parts
of the helper to vary across occurrences, so the same pattern can match
many different concrete subprograms (e.g., the same structure with different
inputs)~\cite{dechter2013bootstrap,ellis2023dreamcoder,cao2023babble,bowers2023top}. This increases flexibility but also enlarges the search
space. Moreover, matching is no longer exact: a helper may apply at many
different locations inside a program, and one must search for all places
where the pattern fits~\cite{bowers2023top}. The result below shows that even before these
additional layers, the core combinatorial problem is already NP-hard.

\par\noindent
\paragraph{Syntax and programs.}
Let $\Sigma$ be a finite set of symbols, and let $f$ be a fixed $n$-ary
constructor. Every program is defined as
\[
p \equiv f(c_1, \dots, c_n),
\]
where each $c_i \in \Sigma$. The arity $n$ is fixed and identical for all
programs.

\par\noindent
\paragraph{Corpus.}
We are given a corpus
\[
\mathcal C \equiv \{p^{(1)}, \dots, p^{(N)}\},
\]
where each program is written as
\[
p^{(j)} \equiv f(c_1^{(j)}, \dots, c_n^{(j)}).
\]

\par\noindent
\paragraph{Helpers (abstractions).}
A reference tuple $(a_1, \dots, a_n) \in \Sigma^n$ is part of the input.
A helper is defined by a subset of positions
\[
S \subseteq \{1, \dots, n\}.
\]

A program $p^{(j)}$ matches $S$ if
\[
\forall i \in S,\quad c_i^{(j)} = a_i.
\]

Let
\[
\mathrm{occ}(S) = \#\{ j : p^{(j)} \text{ matches } S \}, 
\]  the number of corpus programs consistent with the helper.

Define the utility
\[
U(S) = |S| \cdot \mathrm{occ}(S).
\]

\par\noindent
\paragraph{Example.}
Let $\Sigma = \{A, B, X, Y\}$ and consider programs of arity $3$:
\[
\begin{array}{c|ccc}
 & 1 & 2 & 3 \\
\hline
p^{(1)} & A & A & X \\
p^{(2)} & A & A & Y \\
p^{(3)} & A & B & X \\
\end{array}
\]
Fix the reference tuple $(A, A, X)$.

If $S = \{1,2\}$, then $p^{(1)}$ and $p^{(2)}$ match, so
$\mathrm{occ}=2$ and $U=4$.
If $S = \{1\}$, all programs match, so $U=3$.

\par\noindent
\paragraph{Decision problem.}
\begin{quote}
\textsc{Best-Single-Helper} \\
\textbf{Input:} A corpus $\mathcal C$, a reference tuple $(a_1,\dots,a_n)$, and an integer $\omega$. \\
\textbf{Question:} Does there exist $S \subseteq \{1,\dots,n\}$ such that
\[
U(S) \ge \omega \, ?
\]
\end{quote}

\begin{proposition}
\textsc{Best-Single-Helper} is NP-complete.
\end{proposition}

\begin{proof}

\par\noindent
\paragraph{Membership in NP.}
Given a subset $S$, we can compute $\mathrm{occ}(S)$ by scanning all programs
and checking the constraint at each position, which takes time $O(Nn)$.
Thus $U(S)$ can be computed in polynomial time.

\par\noindent
\paragraph{High-level idea of the reduction.}
We reduce from \textsc{Maximum-Edge-Biclique}: given a bipartite graph,
find subsets $S \subseteq L$ and $T \subseteq R$ such that all edges between
them exist, maximizing $|S|\cdot |T|$. The decision version asks whether
there exists such a biclique with $|S|\cdot|T| \ge k$ , which is known to be NP-complete \citep{peeters2003maximum}.

Positions will correspond to vertices in $L$, and programs to vertices in $R$.
Selecting positions enforces that only programs adjacent to all those vertices
can match, yielding a biclique structure.

\par\noindent
\paragraph{Reduction.}
Let $G = (L \cup R, E)$ be a bipartite graph with
\[
L = \{u_1, \dots, u_n\}, \quad R = \{v_1, \dots, v_m\},
\]
and let $k$ be the target.

Set $M \equiv n + 1$.

\par\noindent
\paragraph{Construction.}
We construct a corpus $\mathcal C$ of size $mM$.
Let $\Sigma$ contain symbols $a_1, \dots, a_n$ and fresh symbols
$b_{ij\ell}$.

For each $v_j \in R$ and each $\ell \in \{1,\dots,M\}$, define
\[
p^{(j,\ell)} \equiv f(c_1^{(j,\ell)}, \dots, c_n^{(j,\ell)}),
\]
where
\[
c_i^{(j,\ell)} =
\begin{cases}
a_i & \text{if } (u_i, v_j) \in E, \\
b_{ij\ell} & \text{otherwise}.
\end{cases}
\]

Set $\omega \equiv Mk$.

This construction is polynomial in the size of $G$.

\par\noindent
\paragraph{Example (matrix view of the construction).} Consider a bipartite graph with 

\[ 
L = \{u_1, u_2\}, \quad R = \{v_1, v_2\}, 
\] 

and edges 

\[ 
(u_1,v_1), (u_2,v_1), (u_1,v_2). 
\]

Thus, $(u_2, v_2) \notin E$. Let $M = 3$. The construction builds a corpus of programs of arity $2$. We can represent it as: 

\[ 
\begin{array}{c|cc} & i=1 & i=2 \\ \hline p^{(1,1)} & a_1 & a_2 \\ p^{(1,2)} & a_1 & a_2 \\ p^{(1,3)} & a_1 & a_2 \\ \hline p^{(2,1)} & a_1 & b_{2,2,1} \\ p^{(2,2)} & a_1 & b_{2,2,2} \\ p^{(2,3)} & a_1 & b_{2,2,3} \\ \end{array} 
\]

\par\noindent
\paragraph{Forward direction (good biclique $\Rightarrow$good helper).}
Suppose there exists a biclique $S \subseteq L$, $T \subseteq R$
such that $|S|\cdot|T| \ge k$.
Let $S$ also denote the corresponding set of indices. For any $v_j \in T$ and any $i \in S$, we have $(u_i,v_j)\in E$,
so $c_i^{(j,\ell)} = a_i$ for all $\ell$.
Thus all $M$ programs associated with $v_j$ match $S$.

Therefore

\[
\mathrm{occ}(S) \ge M|T|, \quad
U(S) \ge M|S||T| \ge Mk = \omega.
\]

\par\noindent
\paragraph{Reverse direction (good helper$\Rightarrow$good biclique).}
Suppose there exists $S \subseteq \{1,\dots,n\}$ such that
\[
U(S) \ge Mk.
\]

Define
\[
S' = \{u_i : i \in S\}.
\]

A program $p^{(j,\ell)}$ matches $S$ if and only if for all $i \in S$,
we have $(u_i,v_j)\in E$. Therefore: If $v_j$ is adjacent to all vertices in $S'$, then all $M$ programs
  $p^{(j,\ell)}$ match $S$; Otherwise, if $v_j$ is missing an edge to some $u_i \in S'$, then for that $i$,
  $c_i^{(j,\ell)} \neq a_i$, so none of the programs $p^{(j,\ell)}$ match $S$.

Let
\[
T = \{v_j \in R : (u_i,v_j)\in E \ \forall u_i \in S'\}.
\]
Then
\[
\mathrm{occ}(S) = M|T|.
\]

Hence
\[
U(S) = |S|\cdot \mathrm{occ}(S) = M|S||T| \ge Mk,
\]
which implies $|S||T| \ge k$.
Thus $(S',T)$ forms a biclique of size at least $k$.

\par\noindent
\paragraph{Conclusion.}
We have shown a polynomial-time reduction from
\textsc{Maximum-Edge-Biclique}, and the problem is in NP,
so \textsc{Best-Single-Helper} is NP-complete.
\end{proof}